\subsection{The Mumford target: rigidity, a numerical criterion, and the
bypasses}\label{ssec:mumfordrigid}

The first case of (F2) is now a propagation problem along one compact
Shimura curve: the exceptional classes are algebraic at every CM point
(\Cref{thm:mumfordcm}), and on every member exactly when $B(W\times_{C}W)$
holds (\Cref{cor:mumfordequiv}). This subsection asks whether the
propagation can be obtained without passing through the Lefschetz standard
conjecture, by one of the mechanisms that have closed other cases of the
Hodge conjecture: a degeneration, a deformation into a larger family on which
the class stays of Hodge type, a specialisation to special points, or
semiregularity. It proves that the first two cannot work, that the third
works exactly when the representing cycles have bounded degree, and it
reduces the fourth to one dimension count at one CM point. It does not prove
the classes algebraic at a very general point.

\begin{setupenv}[The classes $\omega_{a}$]\label{setup:omegaa}
Keep \Cref{setup:mumford} and let $W\to C$ be a family as in
\Cref{cor:mumfordlefschetz}, with fibres $X_{t}$. Over $\bar{\QQ}$ let
$\pi_{0},\pi_{12},\pi_{13},\pi_{23}\in H^{2}(X)\otimes H^{2}(X)\subset
H^{4}(X\times X)$ be the classes that correspond, under the isomorphism
$H^{2}\otimes H^{2}\cong\End H^{2}$ given by the form ${\textstyle\bigwedge^{2}}\psi$
on $H^{2}(X)={\textstyle\bigwedge^{2}}V^{\vee}$, to the projections onto $\QQ\theta$,
$U_{12}$, $U_{13}$ and $U_{23}$ of \Cref{prop:mumfordrm}(i), and for
$a=(a_{0},a_{1},a_{2},a_{3})$ put
\[
  \omega_{a}=a_{0}\pi_{0}+a_{1}\pi_{12}+a_{2}\pi_{13}+a_{3}\pi_{23},
  \qquad L_{0}(a)=a_{0}+a_{1}+a_{2}+a_{3}.
\]
The four tensors are invariant under $G$ and under the exchange $\tau$ of the
two factors of $X\times X$, because ${\textstyle\bigwedge^{2}}\psi$ is symmetric; so each
$\omega_{a}$ is a flat section over $C$ of Hodge type at every point. The
argument of \Cref{prop:mumfordrm}(ii) shows that the rational ones are those
with $a_{0}\in\QQ$ and $(a_{1},a_{2},a_{3})=(\sigma_{3}(\alpha),
\sigma_{2}(\alpha),\sigma_{1}(\alpha))$ for some $\alpha\in F$; we call such a
class \emph{exceptional} when $\alpha\notin\QQ$. The products of divisor
classes in $H^{2}\otimes H^{2}$ are the rational $\omega_{a}$ with
$\alpha\in\QQ$, and $\pi_{0}$ is a multiple of
$\mathrm{pr}_{1}^{*}\theta\cup\mathrm{pr}_{2}^{*}\theta$; so adding a rational
multiple of $\pi_{0}$ to an exceptional class changes $L_{0}$ and not the
question of algebraicity. For $c\in C$ write $Y=X_{c}\times X_{c}$, let
$HT^{2}(Y)=H^{0,2}(Y)\oplus H^{1}(T_{Y})\oplus H^{0}({\textstyle\bigwedge^{2}}T_{Y})$ act on
$H^{*}(Y,\CC)$ by contraction as in \Cref{setup:criterion}, of dimension
$28+64+28=120$, and let $\kappa\in H^{1}(T_{Y})$ be the Kodaira--Spencer class
at $c$ of the family $W\times_{C}W\to C$. Put
\[
  e(\omega)=\dim HT^{2}(Y)-\dim\Ann_{HT^{2}(Y)}(\omega)
  =120-\dim\Ann_{HT^{2}(Y)}(\omega).
\]
\end{setupenv}

\begin{lemma}[Linear relations among conjugates]\label{lem:conjugates}
Let $F$ be a cubic number field with embeddings $\sigma_{1},\sigma_{2},
\sigma_{3}$, let $\alpha\in F\setminus\QQ$, and let
$c_{0},c_{1},c_{2},c_{3},a_{0}\in\QQ$. The three conjugates
$\sigma_{i}(\alpha)$ are pairwise distinct and nonzero, and if
$c_{0}a_{0}+\sum_{i}c_{i}\sigma_{i}(\alpha)=0$ then $c_{1}=c_{2}=c_{3}$.
\end{lemma}

\begin{proof}
The minimal polynomial of $\alpha$ has degree three, since $[F:\QQ]$ is prime,
and is separable, so the conjugates are distinct; they are nonzero because
$\alpha\ne0$. Let $\Gamma$ be the Galois group of a normal closure; it
permutes the $\sigma_{i}$ transitively by $\gamma\circ\sigma_{i}=
\sigma_{\gamma(i)}$. Put $\beta=\alpha-\Tr_{F/\QQ}(\alpha)/3$, so that
$\beta\ne0$ has trace zero, and $r=-c_{0}a_{0}-(c_{1}+c_{2}+c_{3})
\Tr_{F/\QQ}(\alpha)/3\in\QQ$; the hypothesis reads $\sum_{i}c_{i}\sigma_{i}
(\beta)=r$. Applying every $\gamma\in\Gamma$ and averaging, each $\sigma_{j}$
occurs $|\Gamma|/3$ times in the place of each $\sigma_{i}$, so
$r=\frac13(\sum_{i}c_{i})\sum_{j}\sigma_{j}(\beta)=0$. The relations
$c\in\QQ^{3}$ with $\sum_{i}c_{i}\sigma_{i}(\beta)=0$ form a $\Gamma$-stable
subspace $N$ containing $(1,1,1)$. The plane $\{\sum c_{i}=0\}$ is an
irreducible representation of $\Gamma$ over $\QQ$, since $\Gamma$ contains
the $3$-cycles, which have no rational eigenvector in that plane; so if $N$
were larger than the line of $(1,1,1)$ it would contain the plane and be all
of $\QQ^{3}$, and $\sigma_{1}(\beta)=0$. Hence $c\in\QQ(1,1,1)$.\footnote{For a
cyclic cubic field $\Gamma$ is the group of $3$-cycles; otherwise it is the
full symmetric group. The argument uses only that $\Gamma$ contains the
$3$-cycles, which holds for every transitive subgroup of the symmetric group
on three letters.}
\end{proof}

\begin{lemma}[One exceptional class gives the others]\label{lem:oneexceptional}
Let $t\in C$. If one rational exceptional class $\omega$ is algebraic on
$X_{t}\times X_{t}$, then every Hodge class in $H^{2}(X_{t})\otimes
H^{2}(X_{t})$ is algebraic; in particular so are the two exceptional classes
of \Cref{prop:mumfordproduct}.
\end{lemma}

\begin{proof}
At a CM point every Hodge class of $X_{t}\times X_{t}$ is algebraic
(\Cref{thm:mumfordcm}); so let $t$ be another point, where the Hodge group is
$G$, and write $X=X_{t}$. For $x\in H^{2}(X)\otimes H^{2}(X)$ let $E_{x}$ be
the endomorphism $y\mapsto x_{*}(\theta^{2}\cup y)$ of $H^{2}(X,\QQ)$, where
$x_{*}$ is the action of $x$ as a correspondence. By Poincar\'e duality and
the hard Lefschetz isomorphism $\theta^{2}\colon H^{2}\to H^{6}$, the map
$x\mapsto E_{x}$ is a $G$-equivariant isomorphism onto $\End H^{2}(X,\QQ)$,
so it carries the Hodge classes onto the Hodge endomorphisms, which are
$\QQ\times F$ by \Cref{prop:mumfordrm}(ii). It carries the products of
divisor classes, which are invariant under $\Sp(V,\psi)$, onto the
$\Sp(V,\psi)$-equivariant endomorphisms, which are $\QQ\times\QQ$, so it
carries $\omega$ to an element $(r,\beta)$ with $\beta\notin\QQ$. The
algebraic Hodge classes form a subspace $A$ containing the products of divisor
classes, and $A$ is closed under composition: if $x_{1},x_{2}$ are algebraic,
then the composite correspondence $x_{1}\circ\Delta_{*}(\theta^{2})\circ
x_{2}$ is algebraic, and so is its K\"unneth component in
$H^{2}\otimes H^{2}$, since the K\"unneth projectors of $X\times X$ are
polynomials in the graphs of $[2]\times1$ and $1\times[2]$, which act on
$H^{i}(X)\otimes H^{j}(X)$ by $2^{i}$ and $2^{j}$ (as in Step 1 of the proof
of \Cref{thm:lefschetzfamily}); that component $x$ has
$E_{x}=E_{x_{1}}E_{x_{2}}$. So the image of $A$ is a subalgebra of
$\QQ\times F$ containing $\QQ\times\QQ$ and $(r,\beta)$, hence
$\QQ\times\QQ[\beta]=\QQ\times F$, because $\beta$ generates the cubic
field $F$.
\end{proof}

\begin{theorem}[Rigidity of the exceptional classes]\label{thm:mumfordrigid}
Let $c\in C$, $Y=X_{c}\times X_{c}$, and let $\omega=\omega_{a}$ be a
rational exceptional class.
\begin{enumerate}[label=\textup{(\roman*)},leftmargin=2.6em,itemsep=2pt]
\item $\Ann_{H^{1}(T_{Y})}(\omega)=\CC\kappa$ and
  $\Ann_{H^{0,2}(Y)}(\omega)=0$.
\item $\Ann_{H^{0}(\bigwedge^{2}T_{Y})}(\omega)=0$ if and only if
  $L_{0}(a)=a_{0}+\Tr_{F/\QQ}(\alpha)\ne0$. So
  $\Ann_{HT^{2}(Y)}(\omega)=\CC\kappa$ and $e(\omega)=119$ whenever
  $L_{0}(a)\ne0$, which can always be arranged by adding a rational multiple
  of $\pi_{0}$.
\item Let $\mathfrak{S}$ be the Siegel space of weight one Hodge structures on
  $H^{1}(Y,\QQ)$ polarised by $\psi\oplus\psi$, of dimension $36$. The germ at
  $Y$ of the locus in $\mathfrak{S}$ where $\omega$ stays of Hodge type is the
  germ of the curve $\{X_{t}\times X_{t}\}$ through $Y$.
\end{enumerate}
\end{theorem}

\begin{proof}
Over $\CC$ the Hodge decomposition at $c$ is
$H^{1,0}(X_{c})=V_{1}\otimes V_{2}\otimes V_{3}^{1,0}$, as at every point of
$C$, since the Hodge cocharacter lies in the third factor. A maximal torus
$S$ of $\SL(V_{1})\times\SL(V_{2})$ therefore preserves the Hodge
decomposition of $H^{1}(Y,\CC)=V\oplus V$ and acts on $HT^{2}(Y)$ and on
$H^{*}(Y,\CC)$, and so does $\tau$. Contraction is equivariant for both, and
$\omega$ is fixed by both, so $x\mapsto x\lrcorner\,\omega$ preserves the
decomposition of $HT^{2}(Y)$ by type ($H^{0,2}$, $H^{1}(T)$ or
$H^{0}(\bigwedge^{2}T)$, which land in the distinct summands $H^{2,4}$,
$H^{1,3}$ and $H^{0,2}$ of $H^{*}(Y)$), by weight of $S$, and by parity under
$\tau$. The annihilator is the sum of the kernels on the $46$ nonzero
blocks. Item (XLIV) of \Cref{sec:verification} writes each block as a matrix
whose entries are linear in $a$ and computes the Gram determinant of its
columns, a polynomial in $a$ that, for real $a$, vanishes exactly where the
columns become dependent. Every Gram determinant is a positive constant times a product of
the following factors, and the one block with a kernel at a generic point is
the $\tau$-even block of weight zero in $H^{1}(T)$, whose kernel is spanned,
for every $a$, by a vector independent of $a$, which is dropped before the
determinant is taken.
\begin{enumerate}[label=\textup{(\alph*)},leftmargin=2.6em,itemsep=1pt]
\item Eleven quadratic forms $P_{1},P_{2},Q_{1},\dots,Q_{4},R,S_{1},\dots,
  S_{4}$, each a positive combination of squares of rational linear forms
  whose common zeros have two of $a_{1},a_{2},a_{3}$ equal or one of them
  zero, for instance
  $2P_{1}=\sum_{i<j}(a_{i}-a_{j})^{2}$,
  $Q_{4}=(3a_{1}-a_{0}-a_{2}-a_{3})^{2}+16(a_{2}-a_{3})^{2}$ and
  $S_{2}=(a_{0}+a_{2})^{2}+3(a_{1}+a_{3})^{2}+8a_{2}^{2}$.
\item Six positive definite forms $Z_{1},\dots,Z_{6}$, each a positive
  combination of the squares of four independent linear forms.
\item The linear forms $a_{2}$, $a_{3}$, $a_{0}-3a_{1}+a_{2}-3a_{3}$,
  $a_{0}-3a_{1}-3a_{2}+a_{3}$, $a_{0}+9a_{1}-3a_{2}-3a_{3}$, whose
  coefficients on $a_{1},a_{2},a_{3}$ are not all equal.
\item The linear form $L_{0}$, which divides the Gram determinants of the
  nine $\tau$-odd blocks of $H^{0}(\bigwedge^{2}T)$ and no other.
\end{enumerate}
The sums of squares are certified in exact arithmetic by item (XLIV) and by
the Lean theorem \texttt{mumford\_rigidity\_sos}. For a rational exceptional
class the coefficients are real, and by \Cref{lem:conjugates} the factors of
(a), (b) and (c) do not vanish: those of (a) would force two conjugates of
$\alpha$ to coincide or one to vanish, those of (b) vanish only at $a=0$,
and those of (c) would give a rational relation among the conjugates with
unequal coefficients. So every block is injective except the one named, and
$L_{0}$ decides the bivector blocks: if $L_{0}(a)\ne0$ they are injective, and
if $L_{0}(a)=0$ the four blocks of dimension one, whose Gram determinants are
$L_{0}^{2}/16$, lie in the annihilator.

The vector of the exceptional block is $\kappa$. The family of Hodge
structures along $C$ is the orbit of the Hodge structure at $c$ under the
real points of the third factor of $G$, so the Kodaira--Spencer class
$\kappa$ is the component of type $(-1,1)$ of an element of
$\mathfrak{sl}(V_{3})$, the lowering operator acting on both copies of $V$.
It kills $\omega$ because $\omega$ is invariant under $G$, it is fixed by
$\tau$ and has weight zero for $S$, and item (XLIV) checks that it spans the
kernel. This proves (i) and (ii).

For (iii), a point of $\mathfrak{S}$ is a Hodge filtration $F^{1}$ on
$H^{1}(Y,\CC)$, and $\omega$ is of Hodge type at it when $\omega\in F^{2}
H^{4}$, a holomorphic condition. For weight one every tangent direction of
$\mathfrak{S}$ at $Y$ is an element $\xi$ of
$\mathfrak{sp}(\psi\oplus\psi)^{-1,1}\subset H^{1}(T_{Y})$, and it moves
$F^{2}H^{4}$ by the derivation $\xi$ of ${\textstyle\bigwedge^{*}}H^{1}$, which maps
$H^{2,2}$ into $H^{1,3}$; so the Zariski tangent space of the Hodge locus is
$\{\xi:\xi\lrcorner\,\omega=0\}$, which is $\CC\kappa$ by (i).\footnote{This is
Griffiths transversality in its simplest case: for a variation of weight one
there is no horizontality condition, so the tangent space of the Hodge locus
of a class of type $(p,p)$ is the kernel of the contraction into
$H^{p-1,p+1}$, see \cite[Chapter~5]{Voi03}. For a class of higher weight on a
general variety only the horizontal directions enter.} The Hodge locus is an
analytic germ containing the smooth curve germ $\{X_{t}\times X_{t}\}$, whose
tangent line is $\CC\kappa$, and its Zariski tangent space is that line; so
the two germs are equal.
\end{proof}

\begin{remark}[What the divisor classes do instead]\label{rem:rigidcompare}
The same computation at classes with $\alpha\in\QQ$ shows what rigidity
excludes. At $a=(3,5,5,5)$, a combination of products of divisor classes, the
tangent space of the Hodge locus in $\mathfrak{S}$ has dimension $10$, and it
contains the tangent space of the diagonal Siegel space
$\{X'\times X'\}$ of dimension $10$, because that class is invariant under
the whole of $\Sp(V,\psi)$; at $a=(1,1,1,1)$ it has dimension $20$. So the
products of divisor classes move with $X'\times X'$ for every $X'$, the
exceptional classes move with $X_{t}\times X_{t}$ for $t\in C$ only, and every
deformation of $Y$ off the curve, however small, makes them cease to be of
Hodge type. \Cref{fig:rigidity} draws the three loci.
\end{remark}

\begin{figure}[tbp]
\centering
  \includegraphics[width=\textwidth]{fig_rigidity}
\caption{\Cref{thm:mumfordrigid}. Left: the Siegel space $\mathfrak{S}$ of
dimension $36$ near $Y=X_{c}\times X_{c}$ and three Hodge loci through it,
labelled by the tangent dimension $T$ at $Y$: the class with equal
coefficients ($T=20$), the products of divisor classes with $a_{1}=a_{2}=a_{3}$
($T=10$, the diagonal Siegel space), and a rational exceptional class ($T=1$),
whose locus is the compact Mumford curve with its CM points (hollow). The
dashed rim is the boundary of $\mathfrak{S}$, which the curve does not reach.
Right: the nine weight blocks of $\mathfrak{sp}^{-1,1}$ under the torus of
the first two factors, with their dimensions and the forms whose product is
the Gram determinant of each; only the block of weight $(0,0)$ has a kernel,
the direction $\kappa$ of the curve. The same holds in the $46$ blocks of
$HT^{2}(Y)$ off the hyperplane $L_{0}=0$. Below: the sums of squares that
make the determinants nonzero at the conjugates of an irrational element of
$F$.}
\label{fig:rigidity}
\end{figure}

\begin{corollary}[No degeneration and no larger family]
\label{cor:nodegeneration}
Let $\omega$ be a rational exceptional class on $X_{c}\times X_{c}$. Every
connected complex analytic family of polarised abelian eightfolds through
$X_{c}\times X_{c}$ on which $\omega$ stays of Hodge type lies, near each
point of the curve, inside the curve $\{X_{t}\times X_{t}:t\in C\}$, which is
compact. In particular no such family degenerates, and on no member of it
other than the CM points of the curve is $\omega$ a combination of products
of divisor classes or of Weil classes of CM fields acting on the member or on
its quotients.
\end{corollary}

\begin{proof}
By \Cref{thm:mumfordrigid}(iii) the Hodge locus of $\omega$ coincides with the
curve in a neighbourhood of every point of the curve, so a connected family
through $Y$ on which $\omega$ stays of Hodge type maps into the curve near
each of its points over the curve; by connectedness and the identity theorem
its image is the curve. The curve is compact because $D$ is a division
algebra, so the arithmetic group of the Shimura curve is cocompact
\cite{Mum69}.\footnote{By \Cref{setup:mumford} the algebra $D$ is definite at
$\sigma_{1}$ and $\sigma_{2}$, so $D\otimes_{F,\sigma_{1}}\RR$ is Hamilton's
quaternions and $D$ is not $M_{2}(F)$: it is a division algebra, split at the
one real place $\sigma_{3}$. The norm one group of an order in it, embedded in
$\SL_{2}(\RR)$ through $\sigma_{3}$, is a Fuchsian group of finite covolume
with no parabolic element: a parabolic $\gamma\ne\pm1$ would make
$\gamma\mp1$ a nonzero nilpotent element of $D$. A Fuchsian group of finite
covolume without parabolic elements has compact quotient.} For the last sentence: a member of the curve other than a CM
point has Hodge group $G$ (proof of \Cref{cor:mumfordequiv}), and on it
$\omega$ lies outside the span of the products of divisor classes and of the
Weil classes of quotients, by \Cref{prop:mumfordproduct}(ii) and (iii). At a
CM point the class is algebraic already (\Cref{thm:mumfordcm}), and the
countably many CM points do not carry algebraicity to the other members
(\Cref{lem:spreading}).
\end{proof}

The same computation controls deformations that are not algebraic at all,
and so decides whether the twistor lines along which Markman moves sheaves
for Weil classes can be used here.

\begin{proposition}[The Hodge locus among all complex tori, and twistor
lines]\label{prop:notwistor}
Let $\omega$ be a rational exceptional class, $Y=X_{c}\times X_{c}$ with
$c\in C$, and $\Lambda=H_{1}(Y,\ZZ)$. Let $\mathcal{T}$ be the space of
complex structures on the real torus $\Lambda_{\RR}/\Lambda$, a complex
manifold of dimension $64$ with tangent space $H^{1}(T_{Y})$ at $Y$, and
$\mathcal{T}_{\omega}\subset\mathcal{T}$ the closed analytic subset on which
$\omega$ is of type $(2,2)$.
\begin{enumerate}[label=\textup{(\roman*)},leftmargin=2.6em,itemsep=2pt]
\item The connected component of $\mathcal{T}_{\omega}$ that contains $Y$ is
  the universal cover $\mathfrak{H}$ of the Mumford curve, the set of
  squares of members of the family, biholomorphic to the upper half plane.
  Every connected complex analytic family of complex tori, algebraic or not,
  through a member of the Mumford family on which $\omega$ stays of Hodge
  type lies in the family.
\item No twistor line through a member of the Mumford family, and no compact
  connected complex curve in $\mathcal{T}$ through it, keeps $\omega$ of
  Hodge type.
\item $\mathcal{T}_{\omega}$ contains twistor lines nevertheless. The algebra
  $D\otimes_{F,\sigma_{1}}\RR$ is Hamilton's quaternions and acts on
  $V_{\RR}$ through the first factor of $G(\RR)$; for a unit quaternion $j$
  with $j^{2}=-1$, the complex structure $j\oplus j$ on $V_{\RR}\oplus
  V_{\RR}$ makes $\omega$ of type $(2,2)$, and these structures form a twistor
  line. On each factor $\psi(x,jy)$ has signature $(4,4)$, and a member whose
  Mumford--Tate group is $G$ is not an abelian variety. These twistor lines
  lie in components of $\mathcal{T}_{\omega}$ other than $\mathfrak{H}$.
\end{enumerate}
\end{proposition}

\begin{proof}
(i) The complex structures on $\Lambda_{\RR}$ for which all the tensors
invariant under $G$ are of Hodge type are those $J$ with $J^{2}=-1$ in the
real points of $G\cdot\mathbb{G}_{m}$, since a reductive group containing the
homotheties is the stabiliser of its invariant tensors. Let $\mathfrak{H}$ be
the connected component of this closed set that contains $Y$; it is closed in
$\mathcal{T}$ and consists of the $J=1\otimes1\otimes j_{3}$ with
$j_{3}\in\SL_{2}(\RR)$, $j_{3}^{2}=-1$, in the half plane that contains the
structure of $Y$, that is of the squares of the members of the Mumford
family. At each of its points the
first order condition for $\omega$ to stay of type $(2,2)$ along
$\xi\in H^{1}(T_{Y})$ is $\xi\lrcorner\,\omega=0$, as in the proof of
\Cref{thm:mumfordrigid}(iii), where no polarisation was used, and by
\Cref{thm:mumfordrigid}(i) the solutions form the line $\CC\kappa$ tangent to
$\mathfrak{H}$. So near each of its points $\mathcal{T}_{\omega}$ is the smooth
curve germ of $\mathfrak{H}$, and $\mathfrak{H}$ is open in
$\mathcal{T}_{\omega}$. Being also closed and connected, it is a connected
component. A connected analytic family on which $\omega$ stays of Hodge type
maps into $\mathcal{T}_{\omega}$, hence into $\mathfrak{H}$ once one of its
points does.

(ii) A twistor line is a nonconstant holomorphic image of $\PP^{1}$ in
$\mathcal{T}$. If $\omega$ stayed of Hodge type along one through a point of
$\mathfrak{H}$, or along any compact connected curve through it, the curve
would lie in $\mathfrak{H}$ by (i), and a holomorphic map from a compact
connected curve to the upper half plane is constant.

(iii) At $\sigma_{1}$ the quaternion algebra $D$ is definite
(\Cref{setup:mumford}), so the first factor of $G(\RR)$ is the group of unit
quaternions, and $V_{\RR}$ restricted to it is a module over
$D\otimes_{F,\sigma_{1}}\RR\cong\mathbb{H}$. The unit quaternions $j$ with
$j^{2}=-1$ form a two-sphere, and the complex structures given by left
multiplication by them on a quaternionic vector space form a twistor line.
Each $j$ lies in $G(\RR)$, and so does the circle $\cos\vartheta+j\sin\vartheta$,
which fixes $\omega$; so $\omega$ is of type $(2,2)$. Item (XLVII) of
\Cref{sec:verification} realises $V_{\RR}$ inside $V_{1}\otimes V_{2}
\otimes V_{3}$ as the fixed space of the real structure that is quaternionic on
the first two factors and real on the third, and finds that $\psi(x,Jy)$ is
definite for the structures of the Mumford family and of signature $(4,4)$ for
$J=j\otimes1\otimes1$. At a member whose Mumford--Tate group is $G$, the Hodge
classes of degree two on $V_{\RR}\oplus V_{\RR}$ are the forms $\psi\otimes s$
with $s$ a symmetric $2\times2$ matrix, by \Cref{prop:mumfordproduct}(ii), and
the signature of $\psi(\cdot,J\cdot)\otimes s$ is $(4(p+q),4(p+q))$ when $s$ has
signature $(p,q)$, so none is definite, none is a polarisation, and the member
is not an abelian variety. The members of the Mumford family are not on these
lines, so by (i) the lines lie in other components.
\end{proof}

Markman's proof for Weil classes on abelian fourfolds moves a sheaf along
twistor lines on which its Chern character stays of Hodge type. By
\Cref{prop:notwistor} no such line passes through the square of a member of
the Mumford family: the twistor lines on which $\omega$ is of Hodge type exist,
but they carry non-algebraic tori whose Hodge cocharacter lies in a definite
factor of $D$, in components of the Hodge locus that the Mumford family never
meets. The Kuga--Satake class of \Cref{thm:mumfordks} behaves differently. The
Kuga--Satake construction applies to every Hodge structure of K3 type on
$(T,q_{\lambda})$, so that class stays of Hodge type over the whole period
domain of $(T,q_{\lambda})$, of dimension seven, and not only along the curve;
it is the one class attached to the target that deformation methods of this
kind are not excluded from.

The same space $T$ also makes the square a holomorphic symplectic variety,
and this decides what the mechanism of van Geemen and Rapagnetta would have
to supply here. They prove the Hodge conjecture for the general abelian
fourfold of Weil type with trivial discriminant by pulling back the second
Chern class of a hyperk\"ahler sixfold along a map from the fourfold, which
gives an algebraic class of codimension two that is not a product of divisor
classes \cite{vGR26}. For a projective hyperk\"ahler manifold $M$ write
$\sigma_{M}$ for its symplectic form, $q_{M}$ for its Beauville--Bogomolov
form and $T(M)$ for the orthogonal complement of $\NS(M)_{\QQ}$ in
$H^{2}(M,\QQ)$ for $q_{M}$.

\begin{proposition}[The square as a holomorphic symplectic
variety]\label{prop:hksquare}
Let $c\in C$ be a point at which the Mumford--Tate group of $X=X_{c}$ is
$G\cdot\mathbb{G}_{m}$, let $Y=X\times X$, and let $\Psi\in\bigwedge^{2}V$ be
the $G$-invariant bivector dual to $\psi$. For $x\in T$ put
$\iota(x)=(x\otimes1)\Psi$ in $V\otimes V=H^{1}(X,\QQ)\otimes H^{1}(X,\QQ)$, a
K\"unneth component of $H^{2}(Y,\QQ)$, and let $\iota_{2}\colon S^{2}T\to
H^{4}(Y,\QQ)$ be the map $xy\mapsto\iota(x)\cup\iota(y)$. Let $C_{i}\in
S^{2}T_{i}$ be the Casimir element of the trace form of
$\mathfrak{sl}(V_{i})$, and for $\nu\in F$ let
$c_{\nu}=\sum_{i}\sigma_{i}(\nu)C_{i}$, a rational element of $S^{2}T$; for
$\nu\ne0$ it is the element dual to the form $q_{1/\nu}$.
\begin{enumerate}[label=\textup{(\roman*)},leftmargin=2.6em,itemsep=2pt]
\item $\iota$ is an embedding of Hodge structures $T(-1)\to H^{2}(Y,\QQ)$, its
  image is the only irreducible sub-Hodge structure of $H^{2}(Y,\QQ)$ with
  $h^{2,0}=1$, and the holomorphic $2$-form $\sigma_{0}$ spanning its part of
  type $(2,0)$ is symplectic.
\item $\iota_{2}(c_{\nu})=\omega_{a}$ with $a_{0}=3\Tr_{F/\QQ}(\nu)$ and
  $\alpha=4\nu-\Tr_{F/\QQ}(\nu)$, in the notation of \Cref{setup:omegaa}. So
  $\iota_{2}(c_{\nu})$ is exceptional if and only if $\nu\notin\QQ$, and every
  rational exceptional class is some $\iota_{2}(c_{\nu})$ plus a rational
  multiple of $\pi_{0}$.
\item Let $M$ be a projective hyperk\"ahler manifold and
  $f\colon Y\dashrightarrow M$ a rational map with $f^{*}\sigma_{M}\ne0$. Then
  $f^{*}$ maps $T(M)$ isomorphically onto $\iota(T)$ and is compatible with
  products of classes in $T(M)$; $f$ is generically finite onto a subvariety
  of dimension eight on which $\sigma_{M}$ is generically nondegenerate;
  $b_{2}(M)\ge10$; and $\dim M\ge10$.
\item Suppose in (iii) that $M=S^{[n]}$ is the Hilbert scheme of $n$ points
  on a projective K3 surface $S$, and
  let $q_{\lambda}$ be the form on $T$ that corresponds to the intersection
  form of $S$ on $T(S)=T(M)$ under $\iota^{-1}f^{*}$
  (\Cref{prop:mumfordrm}(iii)). If $\lambda\notin\QQ$, or if the real
  multiplication of $T(S)$ is induced by algebraic cycles on $S\times S$,
  then the Hodge conjecture holds for $Y$.
\end{enumerate}
\end{proposition}

\begin{proof}
(i) The map $\iota$ is $G$-equivariant because $\Psi$ is invariant, and
injective. Its image lies in $S^{2}V$: for $x$ in the symplectic Lie algebra
of $\psi$ one has $(x\otimes1)\Psi=-(1\otimes x)\Psi$, and the exchange of the
two factors carries $(x\otimes1)\Psi$ to $-(1\otimes x)\Psi$ because $\Psi$
is alternating. An element $g$ of $G\cdot\mathbb{G}_{m}$ multiplies $\Psi$ by
the multiplier $\chi(g)$ of $\psi$, so $(g\otimes g)\iota(x)=
\chi(g)\,\iota(gxg^{-1})$: the map $\iota$ is equivariant for the adjoint
action on $T$ twisted by $\chi$, which is the action defining $T(-1)$. Both
Hodge structures come from the Hodge cocharacter through
$G\cdot\mathbb{G}_{m}$, so $\iota$ is a morphism of Hodge structures from
$T(-1)$. The sub-Hodge structures
of $H^{2}(Y,\QQ)$ are its $G$-stable subspaces, since $\mathbb{G}_{m}$ acts
by scalars. Over $\bar{\QQ}$, item (XLVIII) of \Cref{sec:verification} splits
$H^{2}(Y)=\bigwedge^{2}(V\oplus V)$ by the three Casimir operators into the
isotypic pieces of highest weights $(0,0,0)$, $(2,2,0)$, $(2,0,2)$,
$(0,2,2)$, $(2,2,2)$, $(2,0,0)$, $(0,2,0)$, $(0,0,2)$, of dimensions
$3,27,27,27,27,3,3,3$, whose parts of type $(2,0)$ have dimensions
$0,0,9,9,9,0,0,1$; the last three pieces are $\iota(T_{1})$, $\iota(T_{2})$,
$\iota(T_{3})$. So an irreducible summand of highest weight $(2,0,2)$ or
$(0,2,2)$ contributes $3$ to $h^{2,0}$, one of weight $(2,2,2)$ contributes
$9$, one of weight $(0,0,2)$ contributes $1$, and the others contribute
nothing. A rational sub-Hodge structure $W$ is stable under the Galois group,
which permutes the highest weights as it permutes the factors, transitively.
So if $W\otimes\bar{\QQ}$ has a summand of weight $(2,2,0)$, $(2,0,2)$ or
$(0,2,2)$, it has summands of all three and $h^{2,0}(W)\ge6$; if it has one of
weight $(2,2,2)$, then $h^{2,0}(W)\ge9$. If $h^{2,0}(W)=1$, therefore,
$W\otimes\bar{\QQ}$ contains the piece $\iota(T_{3})$ of weight $(0,0,2)$,
which has multiplicity one, and with it $\iota(T_{1})$ and $\iota(T_{2})$; so
$W\supseteq\iota(T)$, and $W=\iota(T)$ if $W$ is irreducible. The
part of type $(2,0)$ of $\iota(T)$ is spanned by $\sigma_{0}=\iota(E_{3})$,
where $E_{3}\in\mathfrak{sl}(V_{3})\otimes\CC$ maps $V_{3}^{0,1}$ onto
$V_{3}^{1,0}$, and item (XLVIII) finds that $\sigma_{0}$ pairs $H^{1,0}$ of the
first factor with $H^{1,0}$ of the second through an invertible matrix. So
$\sigma_{0}$ has rank eight: it is a holomorphic symplectic form on $Y$.

(ii) Item (XLVIII) computes
$\iota_{2}(C_{1})=3\pi_{0}-\pi_{12}-\pi_{13}+3\pi_{23}$ and the two cyclic
permutations of this identity, in the split model; both sides are built
from $\psi$ and the actions of the three factors, so the identities hold over
$\bar{\QQ}$ in \Cref{setup:mumford}. Summing with the coefficients
$\sigma_{i}(\nu)$ gives $\iota_{2}(c_{\nu})=3\Tr(\nu)\pi_{0}+
\sum\sigma_{k}(4\nu-\Tr(\nu))\pi_{ij}$, the sum over $\{i,j,k\}=\{1,2,3\}$ with
$i<j$, which is $\omega_{a}$ with the stated $a_{0}$ and $\alpha$. Since
$\Tr(\nu)\in\QQ$, the element $\alpha$ is rational exactly when $\nu$ is; and
$\Tr(\alpha)=\Tr(\nu)$, so $\nu=(\alpha+\Tr(\alpha))/4$ recovers $\nu$ from
any $\alpha\in F$. Two rational classes with the same $\alpha$ differ by a
multiple of $\pi_{0}$. The Casimir element $C_{i}$ is dual to the trace form
on $T_{i}$, so $c_{\nu}$ is dual to $\sum_{i}\sigma_{i}(\nu)^{-1}
\operatorname{tr}(x_{i}y_{i})=q_{1/\nu}(x,y)$.

(iii) Resolve the indeterminacy of $f$ by a sequence of blow-ups along
smooth centres, $\pi\colon\tilde{Y}\to Y$ and $\tilde{f}\colon\tilde{Y}\to M$,
and put $f^{*}=\pi_{*}\tilde{f}^{*}$, a morphism of Hodge structures. The
kernel of $\pi_{*}$ on $H^{2}(\tilde{Y},\QQ)$ is spanned by the classes of
the exceptional divisors, so $\alpha\mapsto\tilde{f}^{*}\alpha-
\pi^{*}f^{*}\alpha$ is a morphism of Hodge structures from $H^{2}(M,\QQ)$ to a
space of Hodge classes. The Hodge structure $T(M)$ is irreducible. A
sub-Hodge structure $W\subseteq T(M)$ with $W^{2,0}=0$ consists of rational
classes of type $(1,1)$, which lie in $\NS(M)_{\QQ}\cap T(M)=0$; if
$W^{2,0}\ne0$, then $\sigma_{M}\in W_{\CC}$, the orthogonal complement of $W$
in $H^{2}(M,\QQ)$ is a sub-Hodge structure orthogonal to $\sigma_{M}$ and
$\bar{\sigma}_{M}$, hence of type $(1,1)$ and contained in $\NS(M)_{\QQ}$,
and $W\supseteq T(M)$. It is not spanned by Hodge classes, so that
morphism vanishes on it, and $\tilde{f}^{*}\alpha=\pi^{*}f^{*}\alpha$ for
$\alpha\in T(M)$. Since $\tilde{f}^{*}$ is a ring homomorphism and
$\pi_{*}\pi^{*}$ is the identity, $f^{*}(\alpha\beta)=f^{*}\alpha\cup
f^{*}\beta$ for $\alpha,\beta\in T(M)$. The kernel of $f^{*}$ on $T(M)$ is a
sub-Hodge structure not containing $\sigma_{M}$, hence zero, and the image is
irreducible with $h^{2,0}=1$, hence $\iota(T)$ by (i). So $f^{*}\sigma_{M}$ is
a nonzero multiple of $\sigma_{0}$ and has rank eight at every point of $Y$
where $f$ is defined; the rank of a pulled-back form is at most the dimension
of the image, so $f$ is generically finite onto its image $Z$, of dimension
eight, and $\sigma_{M}$ is generically nondegenerate on $Z$. As $M$ is
projective, $b_{2}(M)=\dim T(M)+\dim\NS(M)_{\QQ}\ge9+1$; this excludes the
generalised Kummer varieties, whose second Betti number is seven
\cite{Bea83}. Finally $\dim M\ge\dim Z=8$, and suppose $\dim M=8$. Then $f$
is dominant, and for $\alpha\in T(M)_{\CC}$ the relation of Fujiki and
Beauville \cite{Bea83,Fuj87} gives
\[
  \int_{Y}(f^{*}\alpha)^{8}=\int_{\tilde{Y}}\tilde{f}^{*}(\alpha^{8})
  =\deg(f)\int_{M}\alpha^{8}=\deg(f)\,c_{M}\,q_{M}(\alpha)^{4},
  \qquad c_{M}>0 .
\]
Write $f^{*}\alpha=\iota(x)$ with $x=x_{1}+x_{2}+x_{3}\in T_{\CC}$ and
$N_{i}=-\det x_{i}$. The form $x\mapsto q_{M}((f^{*})^{-1}\iota(x))$ is
invariant under $G$, so it is $q_{\lambda}(x,x)=
\sum_{i}\sigma_{i}(\lambda)\operatorname{tr}(x_{i}^{2})=
2\sum_{i}\sigma_{i}(\lambda)N_{i}$ for some $\lambda\in F^{\times}$ by
\Cref{prop:mumfordrm}(iii). On the other side, item (XLVIII) finds
$\iota(x)^{8}=8!\det(x)$ times a volume form, and
$\det(x_{1}\otimes1\otimes1+1\otimes x_{2}\otimes1+1\otimes1\otimes
x_{3})=\Delta(N_{1},N_{2},N_{3})^{2}$ with
$\Delta=N_{1}^{2}+N_{2}^{2}+N_{3}^{2}-2N_{1}N_{2}-2N_{1}N_{3}-2N_{2}N_{3}$;
this is also the product of the eight numbers $\pm r_{1}\pm r_{2}\pm r_{3}$,
where $\pm r_{i}$ are the eigenvalues of $x_{i}$ and $r_{i}^{2}=N_{i}$. The
$N_{i}$ are algebraically independent functions on $T_{\CC}$, so
$\Delta^{2}$ would be a constant multiple of the fourth power of the linear
form $\sum_{i}\sigma_{i}(\lambda)N_{i}$. But $\Delta$ is a nondegenerate
ternary quadratic form, hence irreducible, and a contradiction results. So
$\dim M\ge10$.

(iv) By \cite{Bea83} the map $H^{2}(S,\QQ)\to H^{2}(S^{[n]},\QQ)$ induced by
the universal subscheme $Z_{n}\subset S\times S^{[n]}$ is injective, carries
the intersection form to the Beauville--Bogomolov form, and has
$H^{2}(S^{[n]},\QQ)=H^{2}(S,\QQ)\oplus\QQ\delta$ with $\delta$ algebraic, so
$T(M)=T(S)$. Let $\Delta_{T}$ be the K\"unneth component in $T(S)\otimes
T(S)$ of the class of the diagonal of $S$; it is algebraic, being the class
of the diagonal minus $[\mathrm{pt}\times S]$, $[S\times\mathrm{pt}]$ and
$\sum_{r}d_{r}\times d_{r}^{\vee}$ over dual bases of $\NS(S)_{\QQ}$. Its image
under the correspondence $Z_{n}\times Z_{n}$, pulled back along the diagonal
of $M$, is the class $\sum_{a}e_{a}\cup e_{a}^{\vee}$ over dual bases of
$T(S)$ for the intersection form, which is therefore algebraic on $M$, and by
(iii) its image under $f^{*}$ is $\iota_{2}$ of the element of $S^{2}T$ dual
to $q_{\lambda}$, that is $\iota_{2}(c_{1/\lambda})$, an algebraic class on
$Y$. If $\lambda\notin\QQ$ it is exceptional by (ii). If instead the real
multiplication is algebraic, apply it to the first factor of $\Delta_{T}$:
for $b\in F$ the class $\sum_{a}(m_{b}e_{a})\cup e_{a}^{\vee}$ is algebraic,
and its image is $\iota_{2}(c_{b/\lambda})$, which is exceptional for
$b\notin\QQ\lambda$. In both cases \Cref{lem:oneexceptional} makes every
Hodge class in $H^{2}(X)\otimes H^{2}(X)$ algebraic; the Hodge classes in the
other K\"unneth components of $H^{4}(Y,\QQ)$ are products of divisor classes,
as at the end of the proof of \Cref{thm:mumfordks}; and the Hodge classes of
$Y$ form a ring generated in degrees two and four (proof of
\Cref{cor:mumfordequiv}).
\end{proof}

\begin{remark}[What the mechanism of van Geemen and Rapagnetta would
need]\label{rem:hkroute}
By \Cref{prop:hksquare}(ii) the exceptional classes are the classes
$\iota_{2}(c_{\nu})$ with $\nu\notin\QQ$, the images of the elements dual to
the forms $q_{\lambda}$, $\lambda=1/\nu$, on the part $\iota(T)$ of $H^{2}(Y)$
generated by the symplectic class; the form with $\lambda$ rational gives a
product of divisor classes. They are the analogues
for $(Y,\sigma_{0})$ of the dual Beauville--Bogomolov class, whose pull-back is
what \cite{vGR26} uses. Part (iii) says what a transfer of that argument
needs: a projective hyperk\"ahler manifold of dimension at least ten with
transcendental part $T(-1)$, such as $S_{\mu}^{[n]}$ with $n\ge5$ for a
surface $S_{\mu}$ of \Cref{thm:mumfordks}, and a rational map from $Y$ onto a
symplectic subvariety of dimension eight. For $M=S_{\mu}^{[n]}$ the form of
(iv) is $q_{\mu/e^{2}}$ for the $e\in F^{\times}$ with $\iota^{-1}f^{*}=m_{e}$
on $T(S_{\mu})=T$, so $\lambda\notin\QQ$ for every such map as soon as
$\mu\notin\QQ^{\times}F^{\times2}$. No such map is known. It would give an
algebraic cycle on $S_{\mu}\times X\times X$ inducing an isomorphism of
$T(S_{\mu})$ onto $\iota(T)\subset H^{1}(X)\otimes H^{1}(X)$, a correspondence
of the kind that the Kuga--Satake class supplies in the proof of
\Cref{thm:mumfordks}; so this is the Kuga--Satake route in geometric form,
not an independent one. Maps to K3 surfaces and to hyperk\"ahler eightfolds
cannot serve, by (iii).
\end{remark}

\begin{theorem}[The Mumford criterion is a dimension]
\label{thm:mumfordnumerical}
Let $c\in C$, $Y=X_{c}\times X_{c}$, $\omega=\omega_{a}$ a rational
exceptional class with $L_{0}(a)\ne0$, and $E$ a perfect complex on $Y$ with
$\ch(E)=N\omega$, $N\ne0$, and no other component.
\begin{enumerate}[label=\textup{(\roman*)},leftmargin=2.6em,itemsep=2pt]
\item $\dim\Ext^{2}(E,E)\ge119$.
\item If $\dim\Ext^{2}(E,E)=119$, then $\operatorname{ev}_{E}\colon
  HT^{2}(Y)\to\Ext^{2}(E,E)$ is surjective with kernel $\CC\kappa$, and
  $\sigma_{E}$ is injective. Conversely, if $\sigma_{E}$ is injective, then
  $\operatorname{ev}_{E}$ has rank $119$ and kernel $\CC\kappa$, and $E$
  extends to first order along the curve.
\item If some such $E$ at one point $c$ has $\dim\Ext^{2}(E,E)=119$, then
  $\omega_{t}$ is algebraic for every $t\in C$, the Hodge conjecture holds for
  $X_{t}\times X_{t}$ for every $t$, and $B(W\times_{C}W)$ holds.
\item At every CM point $c$ some perfect complex $E$ on $Y$ has
  $\ch(E)=N\omega$ for some $N\ne0$ and no other component.
\end{enumerate}
\end{theorem}

\begin{proof}
(i) and (ii). By the square $\sigma_{E}\circ\operatorname{ev}_{E}=
(\,\cdot\,)\lrcorner\,\ch(E)$ of \Cref{setup:criterion}, $\ker
\operatorname{ev}_{E}\subseteq\Ann_{HT^{2}}(\omega)=\CC\kappa$ by
\Cref{thm:mumfordrigid}(ii), so $\operatorname{ev}_{E}$ has rank at least
$119$, which is (i). If $\dim\Ext^{2}(E,E)=119$, then
$\operatorname{ev}_{E}$ is onto, its kernel has dimension one and lies in
$\CC\kappa$, so equals it; and $\sigma_{E}(\operatorname{ev}_{E}(x))=
N\,x\lrcorner\,\omega$ vanishes only for $x\in\CC\kappa=\ker
\operatorname{ev}_{E}$, so $\sigma_{E}$ is injective on
$\image\operatorname{ev}_{E}=\Ext^{2}(E,E)$. Conversely, if $\sigma_{E}$ is
injective, then $\sigma_{E}(\operatorname{ev}_{E}(\kappa))=N\kappa\lrcorner
\,\omega=0$ forces $\operatorname{ev}_{E}(\kappa)=0$, so the kernel is
$\CC\kappa$ and the rank is $119$; and $\operatorname{ev}_{E}$ restricted to
$H^{1}(T_{Y})$ is the obstruction map \cite{Tod09}, which vanishes at
$\kappa$.

(iii) By (ii) $\sigma_{E}$ is injective. The Chern character of $E$ is
$N\omega$, a flat section of Hodge type over the whole of $C$, so
\Cref{thm:BF} applies to every infinitesimal deformation of $Y$ inside
$W\times_{C}W\to C$. The proof of \Cref{thm:semiregclosure} now applies word
for word, with the stack of perfect complexes on the fibres of $W\times_{C}
W\to C$ \cite{Lie06} in place of that of $\cA\to\cS$, or the analytic
Kuranishi germ if $\Ext^{<0}(E,E)\ne0$: it is smooth over $C$ at $[E]$, its
image contains a nonempty open subset $U\subseteq C$, and for $t\in U$ a
perfect complex $E_{t}$ on $X_{t}\times
X_{t}$ in the same connected component has $\ch(E_{t})=N\omega_{t}$, so
$\omega_{t}$ is algebraic, and by \Cref{lem:oneexceptional} so are both
exceptional classes. $U$ is uncountable, and \Cref{cor:mumfordequiv} gives
the rest.

(iv) By \Cref{thm:mumfordcm} there is $\xi\in\CH^{2}(Y)_{\QQ}$ with class
$\omega$, and the argument of \Cref{prop:ktheory} gives $E$: the Chern
character is an isomorphism $K_{0}(Y)_{\QQ}\to\CH^{*}(Y)_{\QQ}$
\cite[Ex.~15.2.16(b) and Cor.~18.3.2]{Ful98}, so $\xi=\ch(u)$ for a unique
$u$, and $Nu$ is the class of a perfect complex for a suitable $N$.
\end{proof}

Part (iii) is the Mumford analogue of \Cref{thm:p2numerical}(iii): one object,
at one point, and one number, which is $119$ here and $2n(2n-1)$ for a Weil
family. For a Weil family that number is attained by no object
(\Cref{thm:p2false}), because an object attaining it would have to deform to
non-algebraic tori; the classes here are rigid, so that argument does not
apply (\Cref{rem:p2prime}). The objects of (iv) exist, and the only open question is whether one
of them, or another representative of the same class in $K_{0}$, attains the
bound of (i). We have not found one, and we do not claim that one exists. The
rigidity of \Cref{thm:mumfordrigid} is not an obstacle to it: semiregularity
moves an object along the Hodge locus of its Chern character, and the curve
is exactly that locus, so a semiregular object would propagate along the
curve and nowhere else, which is all that is needed.

\begin{corollary}[Bounded data at CM points]\label{cor:mumfordbounded}
Fix a relatively ample line bundle on $W\times_{C}W\to C$, and let $\omega$ be
a rational exceptional class. The following are equivalent.
\begin{enumerate}[label=\textup{(\alph*)},leftmargin=2.6em,itemsep=2pt]
\item $\omega_{t}$ is algebraic for every $t\in C$; equivalently, by
  \Cref{lem:oneexceptional} and \Cref{cor:mumfordequiv}, the Hodge conjecture
  holds for $X_{t}\times X_{t}$ for every $t\in C$.
\item There are a finite set $\Phi$ of pairs $(\underline{P},\underline{q})$,
  each a tuple of Hilbert polynomials with a tuple of rational numbers of the
  same length, and infinitely many points $c\in C$ at each of which
  $\omega_{c}=\sum_{i}q_{i}[Z_{i}]$ for subschemes $Z_{i}\subset X_{c}\times
  X_{c}$ with Hilbert polynomials $P_{i}$, for some $(\underline{P},
  \underline{q})\in\Phi$.
\end{enumerate}
In particular the classes are algebraic on every member as soon as they are
represented at infinitely many CM points by cycles of bounded degree.
\end{corollary}

\begin{proof}
For a pair $(\underline{P},\underline{q})$ let $\Sigma_{\underline{P},
\underline{q}}\subseteq C$ be the set of points at which $\omega$ is
represented with these data. The proof of \Cref{thm:closedlocus} applies
verbatim to $(W\times_{C}W,C,\omega)$, as in \Cref{lem:spreading}, and shows
that each $\Sigma_{\underline{P},\underline{q}}$ is Zariski closed. If (b)
holds, one pair of $\Phi$ serves at infinitely many points, so its
$\Sigma$ is an infinite closed subset of the irreducible curve $C$, hence
$C$, and (a) holds. If (a) holds, $C$ is the union of the countably many
$\Sigma_{\underline{P},\underline{q}}$, and $C$ is uncountable, so one of
them is infinite, hence equal to $C$; then (b) holds with $\Phi$ that one
pair and every point of $C$.
\end{proof}

\Cref{thm:mumfordcm} gives a representation at every CM point, through
Markman's theorem for abelian fourfolds \cite{Mar25a} and the ring of
\Cref{thm:mumfordcm}, with no control of degree, and isogenies between CM
members, which move CM points along their Hecke orbits, do not bound it
either. \Cref{cor:mumfordbounded} says that such a bound is not a convenience
but the whole remaining content: with it the target falls, and without it the
CM points, being countable, say nothing about a very general member
(\Cref{lem:spreading} needs uncountably many points).

\begin{remark}[The bypass mechanisms, one by one]\label{rem:bypassaudit}
\Cref{fig:bypass} collects what is known about each route to the Mumford
target that avoids the Lefschetz standard conjecture.
\begin{enumerate}[label=\textup{(\arabic*)},leftmargin=2.6em,itemsep=2pt]
\item \emph{Degeneration.} Blocked: the base is compact and the Hodge locus is
  the base (\Cref{cor:nodegeneration}), so no family on which $\omega$ stays
  of Hodge type approaches the boundary, and there is no limit mixed Hodge
  structure to specialise to.
\item \emph{Deformation into a larger Hodge locus}, on which $\omega$ would
  become a product of divisor classes or a Weil class and be algebraic by
  an earlier theorem, then be carried back. Blocked by
  \Cref{thm:mumfordrigid}(iii), already to first order.
\item \emph{Hochschild deformations}: twisting $D^{b}(Y)$ by a $B$-field or a
  Poisson bivector, as in the gerby and noncommutative deformations used for
  Weil classes. Blocked: an object with $\ch(E)=N\omega$ extends along no
  direction of $HH^{2}(Y)$ except that of the curve
  (\Cref{cor:mumfordsplit}(i)).
\item \emph{Constructions from line bundles}: cones, tensor products, duals,
  pull-backs and push-forwards along homomorphisms, and Fourier--Mukai
  functors whose kernels are built in the same way. Blocked at every point of
  the curve: everything so built has a Chern character invariant under the
  Lefschetz group, and the exceptional classes are not
  (\Cref{thm:lefschetzclosure}).
\item \emph{Twistor and hyperholomorphic deformation}, the mechanism of
  Markman's proof for Weil classes: move a sheaf along twistor lines on which
  its Chern character stays of Hodge type. Blocked on the square: no twistor
  line through a member of the Mumford family keeps $\omega$ of Hodge type,
  and those that do carry non-algebraic tori, in other components of the
  Hodge locus (\Cref{prop:notwistor}).
\item \emph{Specialisation.} The class is algebraic at every CM point
  (\Cref{thm:mumfordcm}), on every power of $X_{c}$ \cite[Theorem~1.3]{Li26};
  this suffices exactly when the degrees are bounded
  (\Cref{cor:mumfordbounded}), and no bound is known.
\item \emph{Reduction modulo $p$.} At every closed point of every reduction
  of the curve the class is a $\QQ_{\ell}$-combination of algebraic cycles
  (\Cref{prop:modp}), by Li's theorem \cite[Theorem~1.1]{Li26}; this
  suffices, with no lifting step, exactly when the Hilbert data of the cycles
  are bounded on a Zariski dense set of closed points of the arithmetic model
  (\Cref{thm:modpdensity}), and no bound is known.
\item \emph{Semiregularity.} Reduced to one dimension count at one CM point
  (\Cref{thm:mumfordnumerical}); an object attaining it has large odd
  self-extensions (\Cref{thm:mumfordprofile}), is not built from line
  bundles (\Cref{prop:mumforddivisor}) and splits only in a constrained way
  (\Cref{cor:mumfordsplit}); no such object is known.
\item \emph{Kuga--Satake.} Reduced to the Kuga--Satake class of one K3 surface
  of Picard number thirteen (\Cref{thm:mumfordks}), which the known cases do
  not reach (\Cref{rem:mumfordks}). Its geometric form is the mechanism of
  van Geemen and Rapagnetta \cite{vGR26}, pull-back of the dual
  Beauville--Bogomolov class of a hyperk\"ahler manifold: the exceptional
  classes are exactly such twisted dual forms of the part of $H^{2}$ of the
  square generated by its symplectic class, and a rational map onto a
  symplectic subvariety of some $S_{\mu}^{[n]}$, $n\ge5$, would suffice;
  K3 surfaces and hyperk\"ahler eightfolds cannot serve
  (\Cref{prop:hksquare}, \Cref{rem:hkroute}).
\item \emph{Motivated classes.} Every Hodge class on an abelian variety is
  motivated \cite[Th\'eor\`eme~0.6.2]{And96b}, so $\omega_{t}$ is motivated for
  every $t$; a motivated class is algebraic once the Lefschetz involutions of
  the auxiliary varieties are (\Cref{prop:lefmot}), and for the auxiliary
  varieties of Andr\'e's proof that is not known. Along the family the auxiliary variety
  that suffices is $W\times_{C}W$ (\Cref{thm:lefschetzpropagation}), and its
  Lefschetz standard conjecture is equivalent to the target
  (\Cref{cor:mumfordequiv}): the route moves the problem and does not shorten
  it.
\end{enumerate}
Five of the ten are closed off by theorems; two work at special points
only, the CM points and the closed points of the reductions modulo $p$, and
each of them suffices exactly when a degree bound holds; two are reduced to a
single object, one Kuga--Satake class or one perfect complex with
$\dim\Ext^{2}=119$ at one CM point; and one is equivalent to the target. None
of them is carried out here, and the target remains open. The five that are
not closed off are one statement and two stronger inputs, not five
obligations: any one of them would suffice (\Cref{cor:mumfordone},
\Cref{rem:mumfordneeded}).
\end{remark}

\begin{figure}[tbp]
\centering
  \includegraphics[width=0.96\textwidth]{fig_bypass}
\caption{\Cref{rem:bypassaudit}: the ten routes to the Mumford target that
avoid the Lefschetz standard conjecture. Clay, with a crossed dashed spoke: a
route closed off by a theorem of \Cref{ssec:mumfordrigid} or
\Cref{ssec:mumfordobject}. Grass: a route that works at special points only,
at the CM points or at the closed points of the reductions modulo $p$, and
suffices exactly when a degree bound holds (\Cref{cor:mumfordbounded},
\Cref{thm:modpdensity}). Indigo: a route reduced to one open statement, of
the same kind as the target or equivalent to it.}
\label{fig:bypass}
\end{figure}
